% Learning curve-informed RL loop (redesigned, self-contained)
\begin{figure}
\centering
\begin{adjustbox}{max width=\textwidth, center}
\begin{tikzpicture}[
    font=\large,
    box/.style   = {rectangle, rounded corners=8pt, draw=black, line width=1pt,
                    minimum height=2.8cm, text centered, align=center, fill=white},
    policy/.style= {box, fill=blue!10,   draw=blue!55!black,   line width=1.4pt, minimum width=6.6cm},
    env/.style   = {box, fill=gray!12,   draw=gray!55!black,   line width=1.4pt, minimum width=6.6cm},
    deriv/.style = {box, fill=orange!20, draw=orange!80!black, line width=1.8pt, minimum width=6.6cm},
    state/.style = {box, fill=blue!4,    draw=black,           line width=1.4pt, minimum width=33cm, minimum height=2.8cm},
    fwd/.style   = {-{Stealth[length=6mm,width=5mm]}, line width=1.8pt},
    fb/.style    = {-{Stealth[length=6mm,width=5mm]}, line width=1.8pt, gray!55!black},
    fbo/.style   = {-{Stealth[length=6mm,width=5mm]}, line width=2pt, orange!80!black},
    lab/.style   = {font=\normalsize, fill=white, inner sep=3pt},
]

% ---- state row spanning the top ----
\node[state] (state)
     {\textbf{\large State}\;\;
      $s_t=\big(s_{t-1},\;(\lambda_t,\; r_t,\; \textcolor{orange!75!black}{c'(k;\lambda_t)},\; \textcolor{orange!75!black}{c''(k;\lambda_t)})\big)$
      \;\;$\oplus$\;\; static feature $d$};

% ---- middle action row ----
\node[env]    (env)  [below=4.4cm of state] {\textbf{Environment}\\[3pt]{\normalsize train config $\to$ log curve}};
\node[policy] (pi)   [left=7.5cm of env]    {\textbf{Policy} $\pi_\theta$\\[3pt]{\normalsize shared LSTM}};
\node[deriv]  (der)  [right=7.5cm of env]   {\textbf{Derivative}\\\textbf{Extractor}};

% ---- forward flow (left -> right) ----
\draw[fwd] (pi)  -- node[lab,above]{config $\lambda_{t+1}$} (env);
\draw[fwd] (env) -- node[lab,above]{curve $c(k)$}           (der);

% ---- feedback to state (vertical) ----
% policy reads the state (down)
\draw[fb]  (state.south -| pi.north) -- node[lab,left]{read} (pi.north);
% environment returns the reward (up)
\draw[fb]  (env.north) -- node[lab,left]{reward $r_{t+1}$}   (env.north |- state.south);
% derivative features close the loop (up) -- highlighted contribution path
\draw[fbo] (der.north) -- node[lab,right,text=orange!80!black]{\textbf{$c',c''$ features}} (der.north |- state.south);

% ---- legend tag for the contribution ----
\node[font=\normalsize, text=orange!80!black, anchor=west] at ($(der.south)+(-3.3cm,-1.0cm)$)
     {\textbf{$\Leftarrow$ our addition (zero extra cost)}};

\end{tikzpicture}
\end{adjustbox}
\caption{\textbf{Learning curve-informed RL loop.} The policy reads the state and picks
a configuration $\lambda_{t+1}$; the environment trains it and logs the learning curve
$c(k)$; the \textcolor{orange!80!black}{\textbf{Derivative Extractor}} turns that curve
into first/second-order features $c',c''$ appended to the next state --- structured
signal at \emph{zero extra evaluation cost}.}
\label{fig:framework}
\end{figure}
